% ============================================================================
% Section IV: Solution Methods
% ============================================================================

\section{Solution Methods}
\label{sec:methods}

\gtopt{} implements three complementary solution methods, selectable via the
\texttt{method} option.  All three solve the same underlying GTEP formulation
(Section~\ref{sec:formulation}); they differ in how they decompose the
problem across the time and scenario dimensions.

\subsection{Monolithic LP}
\label{sec:monolithic}

The default method assembles the entire multi-scenario, multi-stage problem
into a single sparse LP/MIP and passes it to the solver backend.  The LP
is constructed in compressed sparse column (CSC) format via the
\texttt{FlatLinearProblem} class.

For a system with $N$ buses, $G$ generators, $D$ demands, $L$ lines,
$E$ batteries, $B$ blocks per stage, $T$ stages, and $S$ scenarios, the
approximate problem size is:

\begin{table}[h]
  \centering
  \caption{LP Problem Size Estimates}
  \label{tab:lp_size}
  \begin{tabular}{@{}ll@{}}
    \toprule
    Quantity & Approximate count \\
    \midrule
    Operational variables & $(G{+}2D{+}2L{+}3E{+}N) \times S \times T \times B$ \\
    Investment variables  & $(G{+}D{+}L{+}E) \times T$ \\
    Bus balance rows      & $N \times S \times T \times B$ \\
    Kirchhoff rows        & $L \times S \times T \times B$ (if DC OPF) \\
    SoC balance rows      & $E \times S \times T \times B$ \\
    Capacity rows         & $(G{+}D{+}L{+}E) \times T$ \\
    \bottomrule
  \end{tabular}
\end{table}

The monolithic method is exact (no decomposition error) and is recommended
for problems where the total number of LP variables fits in memory.  For
the IEEE~57-bus system with 24 hourly blocks, the LP has approximately
$57 \times 24 \times (7 + 84 + 160 + 57) \approx 420\,000$ variables.

\subsection{Stochastic Dual Dynamic Programming (SDDP)}
\label{sec:sddp}

For multi-stage stochastic problems with many scenarios, \gtopt{} implements
SDDP following the algorithm of Pereira and
Pinto~\cite{pereira1991sddp}.  The stages are decomposed into
individual subproblems connected by Benders optimality cuts that
approximate the expected future cost.

\begin{algorithm}[h]
  \caption{SDDP Algorithm in \gtopt{}}
  \label{alg:sddp}
  \begin{algorithmic}[1]
    \STATE Initialize: empty cut sets $\mathcal{C}_t = \emptyset$
      for all $t$
    \FOR{iteration $k = 1, 2, \ldots, K$}
      \STATE \textbf{Forward pass}: Solve stages $t=1,\ldots,T$
        sequentially for sampled scenarios; record state variables
        (SoC, reservoir volumes)
      \STATE \textbf{Backward pass}: For $t=T,\ldots,1$, solve
        each stage for all aperture scenarios; compute dual
        variables; generate Benders cuts:
        \[
          \alpha_t \geq \hat{\alpha}_t + \hat{\pi}_t^\top
            (x_t - \hat{x}_t)
        \]
      \STATE Add cuts to $\mathcal{C}_t$
      \STATE Check convergence: upper/lower bound gap $\leq \epsilon$
    \ENDFOR
  \end{algorithmic}
\end{algorithm}

\gtopt{}'s SDDP implementation includes several enhancements:

\begin{itemize}
  \item \textbf{Apertures}: configurable backward-pass scenario sampling
    that groups scenarios into subsets for cut generation, reducing
    computational cost while preserving solution quality.
  \item \textbf{Elastic filters}: infeasible subproblems are handled by
    relaxing state variable bounds with penalty costs, ensuring the
    algorithm does not stall.
  \item \textbf{Boundary cuts}: pre-computed future cost approximations
    at the terminal stage, enabling rolling horizon studies.
  \item \textbf{Hot start}: cuts from previous runs can be loaded to
    accelerate convergence.
  \item \textbf{Cut management}: dominated cuts are periodically purged
    to control LP size growth.
\end{itemize}

\subsection{Cascade Decomposition}
\label{sec:cascade}

The cascade method is a novel multi-level hybrid that decomposes the
problem hierarchically:

\begin{enumerate}
  \item \textbf{Level~0 (strategic)}: A coarsened version of the problem
    (aggregated blocks, reduced scenarios) is solved monolithically to
    determine investment decisions and approximate cost-to-go functions.
  \item \textbf{Level~1 (tactical)}: SDDP is applied to the full
    stochastic problem using the Level~0 investments as warm-start and
    the cost-to-go cuts as initial approximations.
  \item \textbf{Level~2+ (operational)}: Additional refinement levels
    can solve detailed operational subproblems with fixed investments.
\end{enumerate}

The cascade method inherits SDDP's convergence guarantees while using the
monolithic coarse solution to dramatically reduce the number of SDDP
iterations required.  It is particularly effective for long-horizon problems
(10+ stages) with many scenarios.

\subsection{Solver Backend Architecture}

All three methods delegate LP/MIP solution to a pluggable solver backend
loaded as a shared library at runtime.  Currently supported backends are:

\begin{itemize}
  \item \textbf{CLP/CBC} (COIN-OR): Default open-source backend.  CLP for
    LP, CBC for MIP~\cite{forrest2005cbc}.
  \item \textbf{CPLEX} (IBM): Commercial high-performance solver, accessed
    via the C Callable Library.
  \item \textbf{HiGHS}: Open-source solver with parallel simplex and
    interior-point methods~\cite{huangfu2018highs}.
\end{itemize}

The backend is auto-detected by priority (CPLEX $>$ HiGHS $>$ CBC $>$ CLP)
or explicitly selected via the \texttt{--solver} flag.  Solver-specific
options (algorithm, threads, tolerances, time limits) are configured
through the \texttt{solver\_options} JSON section.
